\documentclass{beamer}
\usetheme{default}
\usecolortheme{seahorse}
\usepackage[utf8]{inputenc}
\usepackage{graphicx}
\usepackage{booktabs}
\usepackage{amsmath}
\usepackage{natbib}
\usepackage[english]{babel}
\usepackage{hyphenat}
\usepackage{tikz}
\usetikzlibrary{shapes,arrows.meta,positioning}

\title[Lecture 4]{Latent Variables: Reflective or Formative?}
\subtitle{Advanced Master in Agricultural Economics and Policy}
\author{Giovanbattista Califano}
\date[Portici, 27/04/2026]{Attitudinal Scales for the Study \\ of Consumer Preferences \\April 27, 2026}
\institute[UNINA]{University of Naples Federico II, giovanbattista.califano@unina.it}

\AtBeginSection[]
{
  \begin{frame}<beamer>
    \frametitle{Moving on to...}
    \tableofcontents[currentsection]
  \end{frame}
}

\begin{document}

\frame{\titlepage}

\begin{frame}{Roadmap}
    \begin{center}
        \begin{tabular}{|l|l|c|}
        \hline
        20/04 &  Hello, Psychometrics! & $\checkmark$ \\
        \hline
        23/04 &  The Questionnaire & $\checkmark$ \\
        \hline
        24/04 & Reliability and Validity of a Measure & $\checkmark$ \\
        \hline
        27/04 & Latent Variables: Reflective or Formative? & $\bigcirc$ \\
        \hline
        30/04 & A Bit of SEM & $\bigcirc$ \\
        \hline
        07/05 & Stata Stata Stata Stata Stata Stata Stata & $\bigcirc$ \\
        \hline
        \end{tabular}
    \end{center}
\vspace{0.1cm}
\footnotesize
Recommended readings: \\
\begin{itemize}
    \item Chapters 4 and 7 from \cite{jhangiani2019research}
    \item Chapters 7, 10 and 14 from \cite{olivero2022psicologia}
    \item Chapter 12 from \cite{mehmetoglu2022applied}
\end{itemize}
\scriptsize
\bibliographystyle{apalike}
\bibliography{text}
\end{frame}

\section{Brief Recap}
\begin{frame}{Brief recap}
    We have already seen how we often wish to study concepts that are \textit{not directly observable and measurable}. \\
    \vspace{1em}
    Through \textit{scaling} techniques, however, we can measure the \textit{reflections} (or \textit{manifestations}) of the construct we are interested in studying.
\end{frame}

\begin{frame}{Brief recap}
  \begin{block}{For example}
    People who are more neophobic towards food technologies will tend to agree with these statements (precisely because they are more neophobic!)
  \end{block}
  \vspace{0.3cm}
  \footnotesize
  \begin{tabular}{l l}
    \textbf{FTNS 1:} & I don't need to eat food made using new technology \\
    & because the food I already eat is fine \\
    \textbf{FTNS 2:} & New food technologies are something I am \\
    & uncertain about \\
    \textbf{FTNS 3:} & New food products are not healthier than \\
    & traditional foods \\
    \textbf{(...)}
  \end{tabular}
  \vspace{0.2cm}
  \begin{center}
    \footnotesize
    \begin{tabular}{ccccc}
      \toprule
      \textit{Strongly} & \textit{Disagree} & \textit{Neither agree} & \textit{Agree} & \textit{Strongly} \\
      \textit{disagree} &  & \textit{nor disagree} &  & \textit{agree} \\
      \midrule
      1 & 2 & 3 & 4 & 5 \\
      \bottomrule
    \end{tabular}
  \end{center}
\end{frame}

\begin{frame}{Brief recap... and a problem?}
    The scores assigned to the responses for each item can therefore be aggregated, by sum or mean. The aggregated score represents, in this case, the participant's level of neophobia towards food technologies.\\
    This classical approach is still widely used in the literature, but presents at least one problem... \\
    \pause
    \vspace{1em}
    \textbf{It assumes that the items used for the aggregated score are a perfect measurement of the construct:} \\
    \[
    FTN_i = \frac{\sum_{j=1}^{k} FTNS_{ij}}{k}
    \]
    \small
    In this average of items:
    \begin{itemize}
        \item No error is accounted for
        \item Items have equal weight in the aggregated score
        \item FTN is treated as a \textit{consequence} of the responses to the scale items, but the causal direction should be reversed
    \end{itemize}
\end{frame}

\section{Latent Variables}
\begin{frame}{What we want}
\[
\begin{aligned}
FTNS_1 &= \lambda_1 \cdot FTN + \varepsilon_1 \\
FTNS_2 &= \lambda_2 \cdot FTN + \varepsilon_2 \\
FTNS_3 &= \lambda_3 \cdot FTN + \varepsilon_3 \\
FTNS_4 &= \lambda_4 \cdot FTN + \varepsilon_4 \\
&\vdots \\
FTNS_k &= \lambda_k \cdot FTN + \varepsilon_k
\end{aligned}
\]
\begin{itemize}
  \item The direction of causality is from the construct ($FTN$) to the individual items/indicators $FTNS_j$, $j=1,\dots,k$.
  \item The observed variables $FTNS_j$ are manifestations of the construct, each subject to measurement error $\varepsilon_j$, $j=1,\dots,k$.
  \item The coefficients $\lambda_j$ are the factor loadings that measure the association between the latent factor $FTN$ and item $j$.
\end{itemize}
\end{frame}

\tikzset{
  lv/.style = {draw=blue!80, fill=blue!10, very thick, ellipse, minimum height=1cm, minimum width=2cm, align=center},
  mv/.style = {draw=blue!80, fill=blue!10, very thick, rectangle, minimum height=0.5cm, minimum width=1cm, align=center},
  ev/.style = {draw=blue!80, fill=blue!10, very thick, circle, minimum size=0.6cm, inner sep=1pt, align=center},
  arrow/.style = {thick, ->, >=Stealth}
}

\begin{frame}{In SEM language}
    \centering
    \begin{tikzpicture}[node distance=0.8cm and 2.2cm, auto]
        \node[mv] (i1) {FTNS 1};
        \node[mv, below=of i1] (i2) {FTNS 2};
        \node[mv, below=of i2] (i3) {FTNS 3};
        \node[mv, below=of i3] (i4) {FTNS 4};
        \node[mv, below=of i4] (ik) {FTNS k};
        \node[ev, left=of i1] (e1) {$\varepsilon_1$};
        \node[ev, left=of i2] (e2) {$\varepsilon_2$};
        \node[ev, left=of i3] (e3) {$\varepsilon_3$};
        \node[ev, left=of i4] (e4) {$\varepsilon_4$};
        \node[ev, left=of ik] (ek) {$\varepsilon_k$};
        \node[lv, right=of i3] (FTN) {Food Technology \\ Neophobia};
        \draw[arrow] (FTN) -- (i1);
        \draw[arrow] (FTN) -- (i2);
        \draw[arrow] (FTN) -- (i3);
        \draw[arrow] (FTN) -- (i4);
        \draw[arrow] (FTN) -- (ik);
        \draw[arrow] (e1) -- (i1);
        \draw[arrow] (e2) -- (i2);
        \draw[arrow] (e3) -- (i3);
        \draw[arrow] (e4) -- (i4);
        \draw[arrow] (ek) -- (ik);
        \draw[loosely dotted, very thick] (i4) -- (ik);
        \draw[loosely dotted, very thick] (e4) -- (ek);
    \end{tikzpicture}
\end{frame}

\tikzset{
  lv/.style = {draw=blue!80, fill=blue!10, very thick, ellipse, minimum height=0.6cm, minimum width=1.2cm, align=center},
  mv/.style = {draw=blue!80, fill=blue!10, very thick, rectangle, minimum height=0.4cm, minimum width=0.8cm, align=center},
  arrow/.style = {thick, ->, >=Stealth}
}

\begin{frame}{In SEM language}
  \centering
  \begin{tabular}{@{}c l@{}}
    % Rectangles and squares
    \tikz[baseline=-0.5ex] \node[mv] {};
    &
    Manifest/observed variables \\[1em]
    % Circles and ellipses
    \tikz[baseline=-0.5ex] \node[lv] {};
    &
    Latent/unobserved variables \\[1em]
    % Unidirectional arrows
    \begin{tikzpicture}[baseline=-0.5ex]
      \draw[arrow] (0,0) -- (0.8,0);
    \end{tikzpicture}
    &
    Causal relationships
  \end{tabular}
\end{frame}

\begin{frame}{Latent variables: Reflective vs Formative}
    \begin{block}{Reflective constructs}
        Until now we have spoken of \textbf{reflective} latent variables: they refer to unobservable theoretical concepts (such as attitudes, emotions, etc.) that are assumed to exist independently of the scientific investigation.
    \end{block}
    But not all latent variables are reflective in nature... \\
    \pause
    \begin{block}{Formative constructs}
        \textbf{Formative} constructs \textit{emerge} from reality as composite variables obtained from a linear combination of multiple observed variables. They are never natural phenomena, but artificial constructions created to achieve specific purposes (such as quality indices, satisfaction, etc.)
    \end{block}
\end{frame}

\section{Reflective Constructs}
\begin{frame}{Reflective latent variables}
    \begin{columns}
    \column{0.5\textwidth}
    \[
    \begin{aligned}
    x_1 &= \lambda_1 \cdot \eta + \varepsilon_1 \\
    x_2 &= \lambda_2 \cdot \eta + \varepsilon_2 \\
    x_3 &= \lambda_3 \cdot \eta + \varepsilon_3 \\
    \end{aligned}
    \]
    \column{0.5\textwidth}
    \begin{tikzpicture}
        \node[lv] (l) {$\eta$};
        \node[mv, left=of l] (x2) {$x_2$};
        \node[mv, above=of x2] (x1) {$x_1$};
        \node[mv, below=of x2] (x3) {$x_3$};
        \draw[arrow] (l) -- (x1);
        \draw[arrow] (l) -- (x2);
        \draw[arrow] (l) -- (x3);
    \end{tikzpicture}
    \end{columns}
    \vspace{1em}
    \pause
    \begin{itemize}
        \item Variations in the latent variable $\eta$ manifest as changes in \textbf{every} observed variable $x$.
        \item The direction of causality is from the construct to the indicators.
        \item Each observed variable is treated as a manifestation, affected by a specific error, of the latent variable.
    \end{itemize}
\end{frame}

\begin{frame}{Reflective latent variables}
    \begin{columns}
    \column{0.5\textwidth}
    \[
    \begin{aligned}
    x_1 &= \lambda_1 \cdot \eta + \varepsilon_1 \\
    x_2 &= \lambda_2 \cdot \eta + \varepsilon_2 \\
    x_3 &= \lambda_3 \cdot \eta + \varepsilon_3 \\
    \end{aligned}
    \]
    \column{0.5\textwidth}
    \begin{tikzpicture}
        \node[lv] (l) {$\eta$};
        \node[mv, left=of l] (x2) {$x_2$};
        \node[mv, above=of x2] (x1) {$x_1$};
        \node[mv, below=of x2] (x3) {$x_3$};
        \draw[arrow] (l) -- (x1);
        \draw[arrow] (l) -- (x2);
        \draw[arrow] (l) -- (x3);
    \end{tikzpicture}
    \end{columns}
    \vspace{1em}
    \pause
    \small
    \begin{itemize}
        \item The indicators of a reflective latent variable are \textit{interchangeable} and we expect a high correlation among them, as they should all be good approximations of the latent construct.
        \item We cannot assign a value to the latent variable, but we can study its relationships with other latent or manifest variables.
        \item The measurement errors $\varepsilon$ allow us to account for the portion of variability in the indicators not explained by the latent factor.
    \end{itemize}
\end{frame}

\begin{frame}{Confirmatory Factor Analysis (CFA)}
    \textbf{Confirmatory factor analysis} allows us to verify whether the hypothesised \textit{reflective} factor structure is supported by the observed data, thus enabling researchers to validate hypotheses about the relationships between observed variables and one or more latent factors.\\
    \pause
    \vspace{1em}
    Various indices are used to evaluate the goodness of fit of a CFA, among the most common: \\
    \vspace{0.5em}
    \resizebox{\textwidth}{!}{
    \begin{tabular}{ll}
        \toprule
        \textbf{Fit Index} & \textbf{Desired Value} \\
        \midrule
        $\chi^2$ (Chi-square) & Close to $0$ \\
        CFI (Comparative Fit Index) & $> 0.95$ (good), $> 0.90$ (acceptable) \\
        TLI (Tucker-Lewis Index) & $> 0.95$ (good), $> 0.90$ (acceptable) \\
        RMSEA (Root Mean Square Error of Approximation) & $< 0.05$ (good), $< 0.08$ (acceptable) \\
        SRMR (Standardized Root Mean Square Residual) & $< 0.08$ \\
        \bottomrule
    \end{tabular}
    }
\end{frame}

\section{Formative Constructs}
\begin{frame}{Formative latent variables}
        \begin{columns}
    \column{0.5\textwidth}
    \[
    \eta = \sum_{j=1}^{J} w_{j} \cdot x_{j}
    \]
    \column{0.5\textwidth}
    \begin{tikzpicture}
        \node[lv] (l) {$\eta$};
        \node[mv, left=of l] (x2) {$x_2$};
        \node[mv, above=of x2] (x1) {$x_1$};
        \node[mv, below=of x2] (x3) {$x_3$};
        \draw[arrow] (x1) -- (l);
        \draw[arrow] (x2) -- (l);
        \draw[arrow] (x3) -- (l);
    \end{tikzpicture}
    \end{columns}
    \vspace{1em}
    \pause
    \small
    \begin{itemize}
        \item The relationship between manifest and latent variables is not to be understood as causal, but rather as prescriptive.
        \item These latent variables are \textbf{formed} by the manifest variables, so their definition depends on the indicators used.
        \item Indicators are specific ingredients: they are not interchangeable and we do not expect a high correlation among them: if we change indicators, the very meaning of the formative latent variable changes.
    \end{itemize}
\end{frame}

\begin{frame}{Examples of formative constructs}
\begin{block}{Socioeconomic Status (SES)}
    \begin{enumerate}
        \item Income
        \item Education level
        \item Occupation
    \end{enumerate}
\end{block}
\begin{block}{Hotel service quality}
    \begin{enumerate}
        \item Staff courtesy
        \item Room cleanliness
        \item Breakfast quality
        \item Check-in efficiency
    \end{enumerate}
\end{block}
\small
\begin{itemize}
    \item If the indicators change, the definition of the construct changes.
    \item Indicators are not necessarily intercorrelated.
\end{itemize}
\end{frame}

\end{document}
